\newcommand{\DoNotLoadEpstopdf}{}
\documentclass[11pt]{article}

\usepackage[margin=1in]{geometry}
\usepackage{amsmath,amssymb,amsfonts,amsthm}
\DeclareFontShape{OT1}{cmr}{m}{scit}{<->ssub*cmr/m/sc}{}
\DeclareFontShape{OT1}{cmr}{bx}{sc}{<->ssub*cmr/bx/n}{}
\usepackage{array}
\usepackage{bm}
\usepackage{graphicx}
\usepackage{multirow}
\usepackage{multicol}
\usepackage{nicefrac}
\usepackage{paralist}
\usepackage{enumitem}
\usepackage{verbatim}
\usepackage{thm-restate}
\usepackage[normalem]{ulem}
\usepackage[compact]{titlesec}
\usepackage{tikz}
\usetikzlibrary{arrows}
\usetikzlibrary{arrows.meta}
\usetikzlibrary{backgrounds}
\usetikzlibrary{calc}
\usetikzlibrary{decorations.markings}
\usetikzlibrary{fit}
\usetikzlibrary{patterns}
\usetikzlibrary{positioning}
\usetikzlibrary{shapes}

\usepackage{fancybox}
\newenvironment{fminipage}%
  {\begin{Sbox}\begin{minipage}}%
  {\end{minipage}\end{Sbox}\fbox{\TheSbox}}

\newenvironment{algbox}[0]{%
\vskip 0.2in
\noindent
\begin{fminipage}{6.3in}
}{%
\end{fminipage}
\vskip 0.2in
}

\usepackage{tcolorbox}
\tcbuselibrary{skins,breakable}
\tcbset{enhanced jigsaw}
\usepackage{mdframed}
\usepackage{xcolor}
\usepackage{cite}
\usepackage{url}
\usepackage{xspace}
\usepackage[mathscr]{euscript}
\usepackage{nameref}
\usepackage[
  linktocpage=true,
  pagebackref=false,
  colorlinks=true,
  linkcolor=blue,
  citecolor=blue,
  urlcolor=blue,
  bookmarks,
  bookmarksopen,
  bookmarksnumbered
]{hyperref}
\usepackage[noabbrev,nameinlink]{cleveref}

\newtheorem{theorem}{Theorem}[section]
\newtheorem{example}{Example}[section]
\newtheorem{lemma}{Lemma}[section]
\newtheorem{proposition}[lemma]{Proposition}
\newtheorem{corollary}[lemma]{Corollary}
\newtheorem{claim}[lemma]{Claim}
\newtheorem{fact}[lemma]{Fact}
\newtheorem{conjecture}[lemma]{Conjecture}
\newtheorem{definition}[lemma]{Definition}
\newtheorem{problem}{Problem}
\newtheorem{remark}[lemma]{Remark}
\newtheorem{observation}[lemma]{Observation}
\renewcommand{\theHlemma}{\thesection.\arabic{lemma}}

\newtheorem*{theorem*}{Theorem}
\newtheorem*{claim*}{Claim}
\newtheorem*{proposition*}{Proposition}
\newtheorem*{lemma*}{Lemma}
\newtheorem*{problem*}{Problem}

\allowdisplaybreaks
\setlength{\parskip}{3pt}
\setlength{\emergencystretch}{2em}

\title{Exact False-Name and One-User Coalition Incentives with Linear Miner Revenue}
\author{}
\date{September 1, 2026}

\begin{document}

\pagenumbering{roman}

\maketitle

\begin{abstract}
Let $h\ge1$ and let $D$ be a distribution on nonnegative values whose public
finite description permits exact sampling in expected polynomial time and
deterministic polynomial-time computation of a positive rational median
$m_D$ and the exact rational probabilities $\Pr[V>m_D]$ and $\Pr[V=m_D]$.
In the infinite-capacity ideal-MPC model, we construct a weakly anonymous
randomized transaction fee mechanism with exact nonnegative rational payments
and miner revenue.  Pointwise, an unconfirmed bid pays zero, a confirmed bid
pays at most its report, and miner revenue is at most total payments.  On
finitely encoded bids comparable with $m_D$ in polynomial time, the mechanism
runs in expected time polynomial in the public-description length, $\log h$,
and the total bid-encoding length.  It is ex post false-name user incentive
compatible for every fixed outsider bid vector and every finite randomized
deviation.  For every miner coalition share $\rho\in[0,1]$, it is exact
Bayesian miner incentive compatible with at least $h$ independent truthful
draws from $D$ and exact Bayesian one-user side-contract-proof with at least
$h$ such honest outsiders and one arbitrary-valued colluding user, against
every finite randomized deviation selected before the honest draws and the
mechanism's private coins.  For a truthful bid vector $V\sim D^\ell$ and every
$\ell\ge h$,
\[
\mathbb{E}[\mu]=
\begin{cases}
m_D/2,&1\le h<64,\\[1ex]
h m_D/128,&h\ge64.
\end{cases}
\]
The same piecewise expression is therefore the infimum of expected revenue
over all $\ell\ge h$, and it is always at least $h m_D/128$.
\end{abstract}

\noindent\textbf{Note.} This paper was fully AI generated through a harness created by Richard Peng that discovers open problems and runs them through an automated research pipeline.

\clearpage
\tableofcontents
\clearpage

\pagenumbering{arabic}

\section{Introduction}
\label{sec:introduction}

A transaction fee mechanism (TFM) decides which submitted transactions are
confirmed, what each confirmed user pays, and how much revenue reaches the
miner.  These choices connect congestion and delay in fee markets with the
long-run security role of transaction fees as block subsidies decline
\cite{HubermanLeshnoMoallemi2021Monopoly,CarlstenEtAl2016Instability}.  The
miner, however, is not a trusted auctioneer: it may submit fake transactions
or coordinate with a user.  Modern transaction fee mechanism design therefore
treats user truthfulness, miner incentives, and collusive deviations as
simultaneous design constraints \cite{Roughgarden2024TFM}.

We study these constraints with infinite capacity and ideal multiparty
computation (MPC).  Infinite capacity removes competition for block space but
not identity manipulation.  A strategic user may submit any finite number of
identities, although at most one identity represents its value-bearing
transaction; payments made by the remaining identities are still real costs.
We require false-name user incentive compatibility (UIC) ex post, for every
fixed collection of outsider bids.  Exact Bayesian miner incentive
compatibility (MIC) and one-user side-contract-proofness ($1$-SCP) require that
a finite false-name deviation selected before the honest values and the
mechanism's private coins are observed offer no expected gain.  Here exact
means zero additive incentive slack.  We also ask for pointwise payment and
budget feasibility and a positive truthful-revenue guarantee uniform over
every population size $\ell\ge h$, when all $\ell$ users draw values
independently from a public distribution $D$ and bid truthfully.

This combination of requirements is sensitive to the implementation and
participation assumptions.  In the plain one-shot model, Chung and Shi show
that UIC together with $1$-SCP forces zero miner revenue, even with infinite
block space \cite{ChungShi2023Foundations}.  Shi, Chung, and Wu introduce the
MPC-assisted TFM model and show that, in their universal-participation,
standard-utility formulation, exact simultaneous user, miner, and one-user
coalition incentives still force zero expected revenue when no positive lower
bound on honest participation is known
\cite{shi_et_al:LIPIcs.ITCS.2023.97}.  A known lower bound $h$, together with
a public value distribution, changes this premise.

\paragraph{Known-$h$ MPC-assisted mechanisms.}
Wu, Shi, and Chung study the same broad known-$h$, infinite-block,
MPC-assisted setting with a public $D$ and fake bids
\cite{wu_et_al:LIPIcs.ITCS.2024.98}.  For $h=1$, their parity posted-price
mechanism has exact incentives against unrestricted finite strategic padding
and earns $m_D/2$, where $m_D$ is a median of $D$.  For general $h$, their
threshold construction permits arbitrary strategic bids and earns
$\Theta(hm_D)$, with exponentially small rather than zero Bayesian incentive
slack.  Their exact linear-programming (LP) construction gives expected
revenue $hm_D/4$ when the total number $d$ of coalition bids satisfies
\[
  d\le \frac18\sqrt{\frac{h}{2\log h}}.
\]
We address the exact general-$h$ regime with no a priori numerical bound on a
finite false-name vector.  This comes with a coalition tradeoff: our theorem
protects against one arbitrary-valued colluding user, whereas the approximate
and bounded-bid results above accommodate larger user coalitions in their
respective regimes.  Our small-$h$ branch uses the parity construction; the
distinct ingredient is the large-$h$ mechanism.

\paragraph{Our contribution.}
For every $h\ge1$ and every explicitly described distribution $D$ satisfying
the exact-computation assumptions below, we construct a weakly anonymous
infinite-capacity ideal-MPC mechanism family $\Pi_{h,D}$.  It is implemented by
\hyperref[alg:parity-mechanism]{\textsc{ParityMechanism}} for $1\le h<64$
and by
\hyperref[alg:sybil-potential-mechanism]{\textsc{SybilPotentialMechanism}}
for $h\ge64$.  On finitely encoded bids comparable with $m_D$, these
mechanisms output confirmation bits, exact nonnegative rational payments, and
exact nonnegative rational miner revenue in expected time polynomial in the
public-description length, $\log h$, and the total bid-encoding length.  They
satisfy the pointwise payment and budget constraints and ex post false-name
UIC.  For every miner coalition share $\rho\in[0,1]$, they satisfy exact
Bayesian MIC in the presence of at least $h$ independent truthful users from
$D$ and exact Bayesian $1$-SCP in the presence of at least $h$ such outsiders
and one arbitrary-valued colluding user.  These guarantees cover every
randomized finite-identity deviation chosen before the honest draws and
private coins.  Moreover, for every $\ell\ge h$, on a truthful bid vector of
$\ell$ independent draws from $D$, expected miner revenue is exactly $m_D/2$
when $h<64$ and exactly $hm_D/128$ when $h\ge64$.  The full statement uses the
explicit-distribution and deviation conventions formalized in
\Cref{def:explicit-distribution,def:model}.

\begin{theorem}[Exact false-name and coalition incentives]
\label{thm:main}
Let $h\ge1$ be an integer.  Let $D$ be a publicly described distribution on
nonnegative values such that its finite description permits exact sampling in
expected polynomial time and deterministic polynomial-time computation of a
positive rational median $m=m_D$ and the exact rational numbers
\[
q_{>} = \Pr_{V\sim D}[V>m],
\qquad
q_{=} = \Pr_{V\sim D}[V=m],
\]
all with polynomially bounded bit length.  There is an infinite-capacity
ideal-MPC mechanism family $\Pi_{h,D}$, implemented by
\hyperref[alg:parity-mechanism]{\textsc{ParityMechanism}} when $1\le h<64$
and by
\hyperref[alg:sybil-potential-mechanism]{\textsc{SybilPotentialMechanism}}
when $h\ge64$, with the following properties.
\begin{enumerate}
  \item It is weakly anonymous and, on finitely encoded bids that can be
  compared with $m$ in polynomial time, runs in expected time polynomial in
  the public-description length, $\log h$, and the total bid-encoding length.
  \item On every finite bid vector it outputs confirmation bits
  $x_i\in\{0,1\}$, nonnegative rational payments $p_i$, and nonnegative
  rational miner revenue $\mu$ satisfying, pointwise,
  \[
  x_i=0\Longrightarrow p_i=0,
  \qquad
  x_i=1\Longrightarrow p_i\le b_i,
  \qquad
  \mu\le\sum_i p_i.
  \]
  \item It is ex post false-name user incentive compatible for every fixed
  outsider bid vector and every randomized deviation with finitely many fake
  identities.  For each miner coalition share $\rho\in[0,1]$, it is exact
  Bayesian miner incentive compatible in the presence of at least $h$
  independent truthful users from $D$, and exact Bayesian
  $1$-side-contract-proof in the presence of at least $h$ independent
  truthful outsiders from $D$ and one colluding user of arbitrary value.
  Bayesian deviations may be randomized and use finitely many fake
  identities, but are chosen without observing honest draws and before
  private mechanism coins.
  \item For every integer $\ell\ge h$,
  \[
  \mathbb{E}_{V\sim D^\ell,\,\Pi_{h,D}}[\mu]
  =
  \begin{cases}
  m_D/2,&1\le h<64,\\[1ex]
  h m_D/128,&h\ge64.
  \end{cases}
  \]
  Consequently, the infimum is determined exactly:
  \[
  \inf_{\ell\ge h}
  \mathbb{E}_{V\sim D^\ell,\,\Pi_{h,D}}[\mu]
  =
  \begin{cases}
  m_D/2,&1\le h<64,\\[1ex]
  h m_D/128,&h\ge64,
  \end{cases}.
  \]
  In particular,
  \[
  \inf_{\ell\ge h}
  \mathbb{E}_{V\sim D^\ell,\,\Pi_{h,D}}[\mu]
  \ge \frac{h m_D}{128}.
  \]
\end{enumerate}
\end{theorem}

\paragraph{Ideas of the construction.}
We first randomize at a possible atom of $D$ at its median so that each honest
bid produces an independent fair eligibility bit.  In the small-$h$ regime,
the parity of these bits masks every fixed finite collection of strategic
bits.  For large $h$, miner revenue is supported on a central binomial window
and normalized so that truthful expected revenue is exactly $hm_D/128$ for
every honest count $\ell\ge h$.  The difficulty is that a deviator can change
the submitted length: eligible fake identities incur payments, but ineligible
ones can pad the input for free.  The large-$h$ mechanism therefore gives an
ineligible bid a small free-confirmation probability that never increases
with the submitted length.  The same length-dependent parameter controls the
payment discount for an eligible bid.  Thus the designated user's direct
utility never increases under padding; once the honest-count premise applies,
an additional identity either lowers the relevant free-confirmation chance or
raises the eligible real bid's payment.  This utility loss is a \emph{Sybil
potential} that covers the largest possible revenue gain from padding.  A
binomial comparison lemma and an exact rational-coin implementation make this
argument uniform over every finite deviation without introducing incentive
slack.

\paragraph{Restricted and reformulated positive regimes.}
Other positive-revenue results alter a different part of the model.  Gafni
and Yaish's infinite-capacity Good Toy Auction on the discrete type space
$\{0\}\cup\mathbb N$ satisfies their dominant-strategy incentive
compatibility (DSIC), myopic miner incentive compatibility (MMIC), SCP, and
off-chain-agreement proofness (OCA-proofness) notions
\cite{GafniYaish2024Discrete}.  Its DSIC definition permits one report and is
not the present all-finite false-name guarantee.  Chen, Simchi-Levi, Zhao, and
Zhou study bounded i.i.d. values, finitely many users and slots, Bayesian user
incentives, one-user collusion, and asymptotic revenue approximation, with
fake-transaction concerns handled through additional sequencing and
commitment structure \cite{ChenSimchiLeviZhaoZhou2025}.  Tang and Yao instead
obtain positive results in a proof-of-stake model with modified intertemporal
utilities and further bid-grid and parameter restrictions
\cite{TangYao2023PoSTFM}.  These are restricted or reformulated versions of
the positive-revenue problem, not weaker theorem statements under the present
definitions.

\paragraph{Organization.}
\Cref{sec:preliminaries} specifies the computational and incentive model and
develops exact median balancing.  \Cref{sec:overview} explains the mechanisms
and the incentive accounting.  \Cref{sec:proof-main} proves the parity
guarantees for small $h$, then establishes the parameters, revenue, and exact
incentive properties of the large-$h$ construction.

\section{Preliminaries}
\label{sec:preliminaries}

\paragraph{Notation and randomness.}
For an event $\mathcal A$, let $\mathbf{1}_{\mathcal A}$ denote its indicator.
For $\alpha\in[0,1]$, $\operatorname{Bernoulli}(\alpha)$ denotes a bit that is
one with probability $\alpha$.  For $n\ge0$,
$\operatorname{Bin}(n,1/2)$ denotes the law of a sum of $n$ independent fair
bits, and $D^\ell$ denotes the product law of $\ell$ independent draws from
$D$.  Unless stated otherwise, distinct primitive randomized decisions use
fresh, mutually independent private coins.

\subsection{Input and incentive model}

We specify the computational input and the deviations quantified by the
incentive properties.

\begin{definition}[Explicit distribution]
\label{def:explicit-distribution}
The public description of a distribution $D$ on $\mathbb{R}_{\ge0}$ is a
finite binary string of length $L$.  It
permits exact sampling from $D$ in expected time polynomial in $L$ and,
deterministically in time polynomial in $L$, outputs a positive rational median
$m=m_D$ together with the exact rational values
$q_{>}=\Pr[V>m]$ and $q_{=}=\Pr[V=m]$.  Their numerators and denominators have
bit length polynomial in $L$.  Here a median satisfies
$\Pr[V\le m]\ge1/2$ and $\Pr[V\ge m]\ge1/2$.  The mathematical mechanism
accepts arbitrary nonnegative real bids.  Its computational guarantee applies
to finite bid encodings for which comparison with $m$ is polynomial-time;
nonnegative binary rationals are a canonical case.  Input size includes $L$,
the binary encoding of $h$, and the entire submitted bid vector.
\end{definition}

\begin{definition}[Ideal-MPC false-name model]
\label{def:model}
On $b=(b_1,\ldots,b_N)$, a mechanism outputs confirmations
$x_i\in\{0,1\}$, payments $p_i\ge0$, and total miner revenue $\mu\ge0$.
Every monetary output is an exact rational represented by a binary numerator
and denominator.  Pointwise feasibility means
$p_i=0$ if $x_i=0$, $p_i\le b_i$ if $x_i=1$, and
$\mu\le\sum_i p_i$.  Infinite capacity means that the mechanism is defined
on every finite bid vector and imposes no additional bound on the number of
confirmations.

A strategic user of value $v\ge0$ may submit any finite vector of identities
and, before mechanism coins are drawn, designate at most one as its real
transaction.  Only the designated identity has intrinsic value.  If $I$ is
the set of submitted identities and $j\in I$ is designated, realized user
utility is
\[
v x_j-\sum_{i\in I}p_i;
\]
if no identity is designated, the value term is absent.  For a coalition
controlling a fraction $\rho\in[0,1]$ of the miners, let $I_C$ contain every
identity submitted by any coalition member.  If its user designates
$j\in I_C$, realized coalition utility is
\[
v x_j-\sum_{i\in I_C}p_i+\rho\mu.
\]
The value term is omitted when no identity is designated, including for a
miner-only coalition.  Honest user behavior is one truthful bid designated as
the user's real transaction, honest miners submit no identities, and aborting
has utility zero.  The ideal functionality enforces the outcome rule above.
Accordingly, miner deviations in this model consist of submitting a finite
vector of identities or aborting; they do not include censoring honest inputs,
observing those inputs before submission, or altering the realized outcome.

Ex post false-name user incentive compatibility compares truthful
single-identity participation with every finite false-name strategy for every
fixed vector of outsider bids, taking expectation only over private mechanism
coins.  For Bayesian miner incentive compatibility, at least $h$ honest users
have independent values from $D$.  For Bayesian $1$-side-contract-proofness,
at least $h$ honest outsiders have independent values from $D$, and one
colluding user has an arbitrary value.  A Bayesian deviation may depend on the
colluding user's value but is selected without observing the honest draws and
before private mechanism coins are generated.  Randomized deviations are
allowed.  Bayesian expectations additionally average over the independent
honest draws.  ``Exact'' means that the corresponding incentive inequality has
zero additive slack.  Weak anonymity
means that permuting the bid vector induces the same permutation of the joint
law of confirmations and payments and leaves the law of $\mu$ unchanged.
\end{definition}

\subsection{Exact coins and balanced eligibility}

For a rational $\alpha=a/s\in[0,1]$, with integers $0\le a\le s$ and
$s\ge1$, the desired object is a bit whose success probability is exactly
$\alpha$.  The following rejection sampler implements this object using fair
bits.

\begin{algbox}
\phantomsection\label{alg:rational-coin}
\textbf{\uline{\textsc{RationalCoin}}}: Exactly sample a rational Bernoulli
variable.

\uline{Input}: A rational number $\alpha=a/s\in[0,1]$, represented by binary
integers $0\le a\le s$ and $s\ge1$.

\uline{Output}: A bit $Z$.

\uline{Guarantee}: $\Pr[Z=1]=\alpha$ exactly.
\begin{enumerate}
  \item Set $k=\lceil\log_2 s\rceil$.
  \item Draw $k$ fair bits as a uniform integer
  $R\in\{0,\ldots,2^k-1\}$.  If $R\ge s$, repeat this step.
  \item Output $Z=\mathbf{1}_{\{R<a\}}$.
\end{enumerate}
\end{algbox}

\begin{lemma}[Exact rational-coin sampling]
\label{lem:rational-coin}
\hyperref[alg:rational-coin]{\textsc{RationalCoin}} outputs one with
probability exactly $\alpha$ and uses expected $O(1+\log s)$ bit operations
and fair random bits.
\end{lemma}

\begin{proof}
Conditioned on acceptance, $R$ is uniform on $\{0,\ldots,s-1\}$, exactly $a$
of whose $s$ elements are below $a$.  Thus the output probability is $a/s$.
If $s>1$, then $2^{k-1}<s\le2^k$, so each iteration is accepted with
probability greater than $1/2$; for $s=1$ it is accepted immediately.  The
expected number of iterations is therefore less than two, and an iteration
uses $O(1+\log s)$ random bits and bit operations.
\end{proof}

The sampler permits exact randomization at an atom of the median.  This is the
only place where the exact probabilities in the public description of $D$ are
needed.

Define
\begin{equation}
\tau=
\begin{cases}
\dfrac{1/2-q_{>}}{q_{=}},&q_{=}>0,\\[1ex]
0,&q_{=}=0.
\end{cases}
\label{eq:tau}
\end{equation}
Exact rational arithmetic on the public values $q_{>}$ and $q_{=}$ computes
$\tau$ in polynomial time and gives it polynomial bit length.
For each submitted bid $b$, its eligibility bit $E(b)$ is sampled independently
according to
\begin{equation}
E(b)=
\begin{cases}
0,&b<m,\\
\operatorname{Bernoulli}(\tau),&b=m,\\
1,&b>m.
\end{cases}
\label{eq:eligibility}
\end{equation}
The middle case is implemented by the one-argument call
$\text{\textsc{RationalCoin}}(\tau)$.

\begin{lemma}[Exact median balancing]
\label{lem:balanced-median}
The number $\tau$ lies in $[0,1]$.  If $V\sim D$, then $E(V)$ is a fair bit.
Eligibility bits produced from independent truthful values are mutually
independent.
\end{lemma}

\begin{proof}
The median property gives
$q_{>}\le1/2\le q_{>}+q_{=}$.  If $q_{=}>0$, these inequalities imply
$0\le1/2-q_{>}\le q_{=}$, so $\tau\in[0,1]$ and
\[
\Pr[E(V)=1]=q_{>}+\tau q_{=}=\frac12.
\]
If $q_{=}=0$, the median inequalities force $q_{>}=1/2$, and the same identity
holds.  Independence follows from independence of the values and the fresh
coin used for each identity.
\end{proof}

\section{Overview}
\label{sec:overview}

The construction has a common fair-bit reduction followed by two mechanisms.
The reduction converts every honest value into an independent fair eligibility
bit without assuming that the distribution is continuous.  For $h<64$, a
parity rule then gives the required revenue and, in the presence of an honest
fair bit, makes strategic reports unable to affect its expectation.  For
$h\ge64$, parity no longer gives revenue linear in $h$; instead we normalize
revenue on a central binomial window.  The main
difficulty in this second branch is that an ineligible false-name bid pays
nothing but changes the submitted length, and can therefore change the revenue
distribution.  A decreasing free-confirmation probability acts as a
\emph{Sybil potential}: the possible revenue gain from such padding is matched
by a loss in the colluding user's direct utility.

\paragraph{The common fair-bit reduction.}
Let $m=m_D$ and write $q_{>}=\Pr[V>m]$ and $q_{=}=\Pr[V=m]$.  At an atom of
the median, the deterministic threshold test need not be balanced, so we split
that atom using
\[
\tau=\frac{1/2-q_{>}}{q_{=}}
\quad\text{when }q_{=}>0,
\qquad \tau=0\quad\text{when }q_{=}=0.
\]
The eligibility rule is zero below $m$, one above $m$, and a
$\operatorname{Bernoulli}(\tau)$ bit at $m$.  The median inequalities ensure
$\tau\in[0,1]$, and for an honest draw
\[
\Pr[E(V)=1]=q_{>}+\tau q_{=} = \frac12.
\]
Thus independent truthful bids yield independent fair bits.  The exact public
median data are used here essentially: they let
\hyperref[alg:rational-coin]{\textsc{RationalCoin}} implement the split with
no approximation error; a sampler-only interface does not in general provide
the exact median masses used in this step.  The same threshold structure also
makes the sign of a user's gain from eligibility agree with the sign of
$v-m$.

\paragraph{The small-$h$ branch.}
For $1\le h<64$,
\hyperref[alg:parity-mechanism]{\textsc{ParityMechanism}} takes the public
pair $(h,D)$ and a finite bid vector, forms the eligibility bits, and outputs
\[
x_i=E_i,\qquad p_i=mE_i,\qquad
\mu=m\mathbf{1}_{\{\sum_iE_i\text{ is odd}\}}.
\]
An eligible fake costs $m$, while an ineligible fake has no effect on the
designated identity.  In the Bayesian incentive comparisons, after fixing all
strategic eligibility outcomes and all but one honest outcome, one independent
fair bit remains; it makes the final parity fair.  Consequently the conditional
expected miner revenue is always $m/2$, regardless of any finite strategic
padding.  This observation proves both coalition properties in this regime,
while the elementary payoff $a(v-m)$ for eligibility probability $a$ proves
ex post false-name user truthfulness.  The same argument gives honest revenue
$m/2\ge hm/128$.

\paragraph{The large-$h$ mechanism.}
Fix $h\ge64$ and set $c=h/128$, $T=\lfloor h/4\rfloor$, and $r=9/10$.  For a
submitted length $N\ge h$, let $B_N\sim\operatorname{Bin}(N,1/2)$ and
introduce
\[
\delta_N=36cr^N,\qquad w_N=(1-\delta_N)m,
\qquad P_N=\Pr[T\le B_N\le N-T],
\qquad
A_N=\frac{c}{(1-\delta_N)P_N}.
\]
Below $h$ we put $\delta_N=1/4$ and set miner revenue to zero; this completes
the mechanism on every finite input while preserving the monotonicity of
$\delta_N$.  On input $(h,D,b_1,\ldots,b_N)$,
\hyperref[alg:sybil-potential-mechanism]{\textsc{SybilPotentialMechanism}}
samples $E_i=E(b_i)$ and outputs
\[
x_i=
\begin{cases}
1,&E_i=1,\\
\operatorname{Bernoulli}(\delta_N),&E_i=0,
\end{cases}
\qquad
p_i=w_NE_i,
\qquad
\mu=w_NA_N\mathbf{1}_{\{T\le K\le N-T\}},
\quad K=\sum_iE_i,
\]
with the displayed revenue rule used when $N\ge h$.  Thus eligible bids pay a
discounted threshold price, whereas ineligible bids receive only the small
free-confirmation probability $\delta_N$.  If all $N$ bids are honest, then
$K\sim\operatorname{Bin}(N,1/2)$, and the choice of $A_N$ gives the exact
normalization
\[
\mathbb{E}[\mu]=w_NA_NP_N=cm=\frac{hm}{128},
\]
independently of $N$.  The central window simultaneously enforces budget
feasibility: a binomial tail estimate yields $A_N\le h/48\le T$, so positive
revenue implies at least $T\ge A_N$ paying bids.

\paragraph{Padding control and the Sybil potential.}
Two quantitative ingredients make the normalization robust to false names.
First, with $Q_N=1-P_N$ and
\[
d_N=c\frac{Q_N}{P_N},
\]
the tail estimate $Q_N\le2e^{-N/8}\le2r^N$ gives
\begin{equation}
\delta_n-\delta_N\ge d_N\quad(h\le n<N),
\qquad
w_NA_N=cm+md_N\le cm+w_N.
\label{eq:overview-potential}
\end{equation}
The first inequality says that increasing the submitted length causes enough
potential drop to cover the largest possible revenue excess $md_N$; the
second says that one eligible fake's payment already covers that excess.

Second, the binomial padding lemma isolates what happens when all added
identities are ineligible.  For $B_L\sim\operatorname{Bin}(L,1/2)$, define
\[
F_{L,N}=\Pr[T\le B_L\le N-T].
\]
Binomial symmetry and unimodality show $F_{L,N}\le P_N$.  Moreover,
$F_{N-1,N}=P_N$, and the same equality holds after adding one fixed eligible
bit.  Hence, in the presence of at least $h$ honest fair bits, arbitrarily many
fixed ineligible identities cannot raise expected revenue above $cm$.  When
$N-1\ge h$ of the bits are honest and fair, the eligibility status of one
additional designated identity leaves expected revenue exactly $cm$.  This
latter equality supplies the honest coalition baseline needed for exact
$1$-side-contract-proofness.

\paragraph{Incentive accounting.}
At total length $N$, if a report $b$ has eligibility probability $a(b)$, its
expected direct utility is
\begin{equation}
u_N(v,b)=\delta_Nv+a(b)(1-\delta_N)(v-m).
\label{eq:overview-direct-utility}
\end{equation}
The coefficient of $a(b)$ changes sign at $m$, so truthful reporting maximizes
this expression.  Its truthful optimum is
\[
U_N(v)=
\begin{cases}
\delta_Nv,&v\le m,\\
v-m+\delta_Nm,&v\ge m,
\end{cases}
\]
which is nonincreasing in $N$.  Adding false names therefore cannot improve a
user's direct payoff for any fixed outsider bid vector, and eligible fakes only
add payments; this proves ex post false-name UIC.

For MIC, the strategic vector is chosen before the honest draws, so the proof
may condition on all fake eligibility outcomes while retaining the honest fair
bits.  If no fake is eligible, the padding lemma bounds expected revenue by
$cm$; if some fake is eligible, its payment $w_N$, together with
\eqref{eq:overview-potential}, covers the entire possible revenue excess even
for a miner share $\rho=1$.  The same estimates handle most of the $1$-SCP
comparison.  Its only additional case is an eligible designated identity
padded solely by ineligible fakes.  If $n$ is the truthful length and $N>n$ is
the padded length, then
\[
U_n(v)-(v-w_N)
\ge m(\delta_n-\delta_N)
\ge md_N.
\]
This loss of user utility covers the same $md_N$ revenue excess, and hence also
its $\rho$-fraction.  These cases exhaust every finite deviation; averaging the
conditional bounds also covers randomized strategies and yields the exact
Bayesian incentive inequalities.

Finally, both mechanisms are symmetric in the submitted bids and hence weakly
anonymous.  Eligibility implies that every positive payment is at most the
corresponding bid; odd parity in the small branch and the central-window
threshold in the large branch bound revenue by total payments pointwise.  The
mechanisms use $O(N)$ calls to
\hyperref[alg:rational-coin]{\textsc{RationalCoin}} and exact rational
arithmetic; in the large branch, $P_N$ is computed from an $O(N)$-term binomial
sum.  This gives exact rational outputs and expected running time polynomial in
the public-description length, $\log h$, and the full bid-encoding length.
Together with the incentive arguments, this proves every guarantee of
\Cref{thm:main}.  In particular, for every honest length $\ell\ge h$,
\[
\mathbb{E}_{V\sim D^\ell,\,\Pi_{h,D}}[\mu]
=
\begin{cases}
m_D/2,&1\le h<64,\\[1ex]
h m_D/128,&h\ge64,
\end{cases}
\]
and hence the infimum over $\ell\ge h$ is at least $hm_D/128$.

\section{Proof of the Main Theorem}
\label{sec:proof-main}

\subsection{The parity mechanism for small \texorpdfstring{$h$}{h}}

The first branch needs no quantitative tail estimates.  Fix $1\le h<64$ and
a finite bid vector $b=(b_1,\ldots,b_N)$.  Mathematically, sample independent
$E_i=E(b_i)$, put $K=\sum_i E_i$, and define
\begin{equation}
x_i=E_i,
\qquad
p_i=mE_i,
\qquad
\mu=m\mathbf{1}_{\{K\text{ is odd}\}}.
\label{eq:parity-rule}
\end{equation}
The algorithmic presentation simply implements these already-defined random
variables.

\begin{algbox}
\phantomsection\label{alg:parity-mechanism}
\textbf{\uline{\textsc{ParityMechanism}}}: Implement the parity rule for the
small-$h$ regime.

\uline{Input}: An integer $1\le h<64$, the public description of $D$, and a
finite bid vector $b=(b_1,\ldots,b_N)$.

\uline{Output}: Confirmation bits $(x_i)_{i=1}^N$, payments $(p_i)_{i=1}^N$,
and miner revenue $\mu$.

\uline{Guarantee}: The outputs obey \eqref{eq:parity-rule}; their incentive,
feasibility, revenue, and resource guarantees are stated in
\Cref{lem:parity-guarantees}.
\begin{enumerate}
  \item Independently sample each $E_i$ according to
  \eqref{eq:eligibility}, using
  $\text{\textsc{RationalCoin}}(\tau)$ when $b_i=m$.
  \item Set $x_i=E_i$ and $p_i=mE_i$ for every $i$.
  \item Compute $K=\sum_{i=1}^N E_i$ and output
  $\mu=m\mathbf{1}_{\{K\text{ is odd}\}}$.
\end{enumerate}
\end{algbox}

\begin{lemma}[Small-regime parity guarantees]
\label{lem:parity-guarantees}
For $1\le h<64$, \hyperref[alg:parity-mechanism]{\textsc{ParityMechanism}}
is weakly anonymous, pointwise feasible, and expected-polynomial-time.  In the
model of \Cref{def:model}, it is ex post false-name user incentive compatible,
exact Bayesian miner incentive compatible, and exact Bayesian
$1$-side-contract-proof for every $\rho\in[0,1]$ and every finite false-name
strategy.  If $V\sim D^\ell$ bids truthfully for any $\ell\ge h$, then
$\mathbb{E}[\mu]=m/2$.
\end{lemma}

\begin{proof}
An eligible bid satisfies $b_i\ge m$, is confirmed, and pays $m$; every other
bid is unconfirmed and pays zero.  If $\mu=m$, oddness implies $K\ge1$, and
hence
\[
\mu=m\le Km=\sum_i p_i.
\]
Thus all pointwise constraints hold.  The rule is permutation-equivariant.
It makes $O(N)$ eligibility decisions, so \Cref{lem:rational-coin} gives
expected polynomial running time with exact rational outputs.

Fix outsider bids and a user of value $v$.  If a report has eligibility
probability $a$, its expected direct utility is $a(v-m)$.  This is maximized by
eligibility probability zero when $v<m$, by probability one when $v>m$, and
by every probability when $v=m$; truthful reporting attains the maximum.
Every eligible fake costs $m$ and every ineligible fake has no effect on the
designated identity.  A strategy with no designated identity has utility at
most zero, while truthful utility is nonnegative.  This proves ex post
false-name user incentive compatibility.

For either Bayesian coalition property, condition on the strategic bid vector,
all of its eligibility outcomes, and all but one honest eligibility bit.  At
least one honest fair bit remains, so \Cref{lem:balanced-median} makes the
parity of $K$ conditionally fair.  Conditional expected miner revenue is
therefore always $m/2$, regardless of finite padding.  In a miner-only
deviation, the miner share remains $\rho m/2$ and fake payments are
nonnegative.  In a one-user coalition, the designated identity's conditional
direct payoff is either zero or $v-m$, and hence is at most the truthful
payoff $\max\{0,v-m\}$; fake payments again can only decrease utility.  The
same conclusions hold after averaging over randomized deviations.  Aborting
gives zero and cannot improve either nonnegative honest baseline.  Finally,
when values are drawn independently from $D$ and bid truthfully, there is an
honest fair eligibility bit, so the same parity argument yields expected
revenue $m/2$.
\end{proof}

\subsection{Parameters and binomial padding for large \texorpdfstring{$h$}{h}}

We now turn to $h\ge64$.  The parameters below normalize honest revenue while
making the free-confirmation probability decrease with the submitted length.
Set
\begin{equation}
c=\frac{h}{128},
\qquad
T=\left\lfloor\frac h4\right\rfloor,
\qquad
r=\frac9{10}.
\label{eq:large-base-parameters}
\end{equation}
For every integer $N\ge0$, define
\begin{equation}
\delta_N=
\begin{cases}
1/4,&N<h,\\
36cr^N,&N\ge h,
\end{cases}
\qquad
w_N=(1-\delta_N)m.
\label{eq:potential-parameters}
\end{equation}
For $N\ge h$, let $B_N\sim\operatorname{Bin}(N,1/2)$ and put
\begin{equation}
P_N=\Pr[T\le B_N\le N-T],
\qquad
Q_N=1-P_N,
\label{eq:central-probability}
\end{equation}
\begin{equation}
A_N=\frac{c}{(1-\delta_N)P_N},
\qquad
d_N=c\frac{Q_N}{P_N}.
\label{eq:normalization-parameters}
\end{equation}
The central interval is nonempty and has positive binomial mass because
$T\le N/4$.

\begin{lemma}[Sybil-potential parameter bounds]
\label{lem:parameter-bounds}
For $h\ge64$, the following properties hold.
\begin{enumerate}
  \item The sequence $(\delta_N)_{N\ge0}$ is nonincreasing and
  $0<\delta_N\le1/4$.
  \item For every $N\ge h$,
  \[
  Q_N\le2e^{-N/8}\le2r^N,
  \qquad
  P_N\ge\frac12.
  \]
  \item For every $N\ge h$,
  \[
  A_N\le\frac h{48}\le T,
  \qquad
  d_N\le4cr^N=\frac{\delta_N}{9}<1-\delta_N.
  \]
  \item If $h\le n<N$, then
  \begin{equation}
  \delta_n-\delta_N\ge d_N.
  \label{eq:potential-drop}
  \end{equation}
  Moreover,
  \begin{equation}
  w_NA_N=\frac{cm}{P_N}=cm+md_N\le cm+w_N.
  \label{eq:maximum-revenue}
  \end{equation}
\end{enumerate}
\end{lemma}

\begin{proof}
Write $B_N-N/2=\sum_{i=1}^N Y_i$, where the independent variables $Y_i$ are
uniform on $\{-1/2,1/2\}$.  For $\lambda\ge0$,
\[
\mathbb{E}[e^{-\lambda Y_i}]=\cosh(\lambda/2)
\le e^{\lambda^2/8}.
\]
Indeed, $\log\cosh z\le z^2/2$ for $z\ge0$ because its derivative satisfies
$\tanh z\le z$ and both sides vanish at zero.  Markov's inequality with
$\lambda=1$ gives
\[
\Pr[B_N-N/2\le-N/4]
\le e^{-N/4}\prod_{i=1}^N\mathbb{E}[e^{-Y_i}]
\le e^{-N/8}.
\]
The upper tail has the same bound by symmetry.  Since $T\le N/4$, the event
outside $[T,N-T]$ is contained in these two tails, proving
$Q_N\le2e^{-N/8}$.  Also $e^{1/8}>10/9$, so $e^{-N/8}\le r^N$.  Finally,
$r^4<2/3$ implies $r^{16}<1/4$ and hence $r^{64}<1/256$.  Thus
$Q_N\le2r^N\le2r^{64}<1/2$, which gives $P_N\ge1/2$.

For every integer $s\ge64$,
\[
\frac{(s+1)r^{s+1}}{sr^s}=r\left(1+\frac1s\right)<1,
\]
so $sr^s$ is decreasing.  Therefore, for $N\ge h$,
\[
\delta_N=36\frac h{128}r^N
\le36\frac h{128}r^h
\le18r^{64}<\frac14.
\]
Below $h$ the sequence equals $1/4$, while above $h$ it decreases by a factor
$r$; this proves the first assertion, including the boundary at $h$.

The bounds $1-\delta_N\ge3/4$ and $P_N\ge1/2$ imply
\[
A_N\le\frac{c}{(3/4)(1/2)}=\frac h{48}.
\]
Since $T=\lfloor h/4\rfloor\ge h/8$ for $h\ge64$, this is at most $T$.
Moreover,
\[
d_N=c\frac{Q_N}{P_N}\le4cr^N=\frac{\delta_N}{9}
<1-\delta_N.
\]
For $N\ge h+1$,
\[
\delta_{N-1}-\delta_N
=36cr^{N-1}(1-r)=4cr^N.
\]
If $h\le n<N$, monotonicity now yields
$\delta_n-\delta_N\ge4cr^N\ge d_N$, proving
\eqref{eq:potential-drop}.  Finally,
\[
w_NA_N=(1-\delta_N)m\frac{c}{(1-\delta_N)P_N}
=\frac{cm}{P_N}=cm+mc\frac{Q_N}{P_N}=cm+md_N.
\]
The inequality $md_N<w_N$ gives \eqref{eq:maximum-revenue}.
\end{proof}

The parameter estimates control the largest possible revenue value.  To
control its expectation after costless ineligible identities are appended, we
next compare central-window probabilities with different numbers of fair
bits.

For fixed $N\ge h$ and $0\le L\le N$, let
$B_L\sim\operatorname{Bin}(L,1/2)$ and define
\begin{equation}
F_{L,N}=\Pr[T\le B_L\le N-T].
\label{eq:padding-probability}
\end{equation}

\begin{lemma}[Binomial padding comparison]
\label{lem:binomial-padding}
For every $N\ge h$ and $0\le L\le N$,
$F_{L,N}\le P_N$, with $F_{N-1,N}=P_N$.  In addition,
\[
\Pr[T\le B_{N-1}+1\le N-T]=P_N.
\]
\end{lemma}

\begin{proof}
Adding one independent fair bit gives
\begin{equation}
F_{L+1,N}-F_{L,N}
=\frac12\bigl(\Pr[B_L=T-1]-\Pr[B_L=N-T]\bigr).
\label{eq:padding-increment}
\end{equation}
If $N-T>L$, the second probability vanishes.  Otherwise put $a=T-1$ and
$b=N-T$.  Then $a\le b$ and $a+b=N-1\ge L$.  If $a\le L/2$, then
$L-b\le a\le L/2$; symmetry and monotonicity of the binomial coefficients up
to their center give
$\binom{L}{a}\ge\binom{L}{L-b}=\binom{L}{b}$.  If $a>L/2$, monotonicity after
the center gives the same inequality directly.  This monotonicity follows
from
\[
\frac{\binom{L}{k+1}}{\binom{L}{k}}=\frac{L-k}{k+1}.
\]
Thus \eqref{eq:padding-increment} is nonnegative.

At $L=N-1$, the indices $T-1$ and $N-T$ sum to $L$, so the two probabilities
in \eqref{eq:padding-increment} are equal.  Hence $F_{L,N}$ is nondecreasing
in $L$ and
$F_{N-1,N}=F_{N,N}=P_N$.  Finally, reflection
$k\mapsto N-1-k$ maps the interval $[T-1,N-T-1]$ onto $[T,N-T]$.
The symmetry of $B_{N-1}$ therefore gives
\[
\Pr[T\le B_{N-1}+1\le N-T]
=\Pr[T-1\le B_{N-1}\le N-T-1]
=F_{N-1,N}=P_N.
\]
\end{proof}

\subsection{The large-\texorpdfstring{$h$}{h} mechanism and its revenue}

The preceding objects now define the large-regime mechanism.  On a bid vector
$b=(b_1,\ldots,b_N)$, sample independent eligibility bits
$E_i=E(b_i)$.  Independently of them, sample
$Z_i\sim\operatorname{Bernoulli}(\delta_N)$ for every $i$, ignoring $Z_i$
when $E_i=1$.  Put
\begin{equation}
x_i=
\begin{cases}
1,&E_i=1,\\
Z_i,&E_i=0,
\end{cases}
\qquad
p_i=w_NE_i,
\qquad
K=\sum_{i=1}^N E_i,
\label{eq:large-user-rule}
\end{equation}
and define
\begin{equation}
\mu=
\begin{cases}
0,&N<h,\\
w_NA_N\mathbf{1}_{\{T\le K\le N-T\}},&N\ge h.
\end{cases}
\label{eq:large-revenue-rule}
\end{equation}
Conditional on the bid vector, all primitive eligibility and free-confirmation
coin draws are mutually independent; the displayed output variables are
functions of those draws and the reports.

\begin{algbox}
\phantomsection\label{alg:sybil-potential-mechanism}
\textbf{\uline{\textsc{SybilPotentialMechanism}}}: Implement the decreasing
free-confirmation and central-window rules.

\uline{Input}: An integer $h\ge64$, the public description of $D$, and a
finite bid vector $b=(b_1,\ldots,b_N)$.

\uline{Output}: Confirmation bits $(x_i)_{i=1}^N$, payments $(p_i)_{i=1}^N$,
and miner revenue $\mu$.

\uline{Guarantee}: The outputs obey
\eqref{eq:large-user-rule}--\eqref{eq:large-revenue-rule}; their feasibility,
anonymity, and resource guarantees are stated in
\Cref{lem:large-implementation}.
{\raggedright
\begin{enumerate}
  \item Compute $\delta_N$ and $w_N$.  Independently sample each $E_i$
  according to \eqref{eq:eligibility}, using
  $\text{\textsc{RationalCoin}}(\tau)$ when $b_i=m$.
  \item For every $i$, if $E_i=1$, set $(x_i,p_i)=(1,w_N)$.  If $E_i=0$,
  set $x_i\leftarrow\text{\textsc{RationalCoin}}(\delta_N)$ and $p_i=0$.
  \item Compute $K=\sum_iE_i$.  If $N<h$, output $\mu=0$; otherwise compute
  $P_N$ and $A_N$ from \eqref{eq:central-probability}--
  \eqref{eq:normalization-parameters} and output
  $\mu=w_NA_N\mathbf{1}_{\{T\le K\le N-T\}}$.
\end{enumerate}
\par}
\end{algbox}

\begin{lemma}[Large-regime feasibility and exact implementation]
\label{lem:large-implementation}
For $h\ge64$,
\hyperref[alg:sybil-potential-mechanism]{\textsc{SybilPotentialMechanism}}
implements the random variables in
\eqref{eq:large-user-rule}--\eqref{eq:large-revenue-rule}, is weakly anonymous,
satisfies every pointwise feasibility constraint in \Cref{def:model}, produces
exact rational monetary outputs, and runs in expected polynomial time.
\end{lemma}

\begin{proof}
Steps 1--3 implement the eligibility, free-confirmation, payment, and revenue
rules in \eqref{eq:large-user-rule}--\eqref{eq:large-revenue-rule} directly.
An eligible bid is at least $m$, is confirmed, and pays
$w_N=(1-\delta_N)m\le m$.  An ineligible bid pays zero whether or not its free
confirmation succeeds, and an unconfirmed bid always pays zero.  If $\mu>0$,
then $N\ge h$, $K\ge T$, and $A_N\le T$ by
\Cref{lem:parameter-bounds}.  Exactly the $K$ eligible bids pay $w_N$, so
\[
\mu=w_NA_N\le w_NT\le w_NK=\sum_i p_i.
\]
If $\mu=0$, budget feasibility is immediate.  Every rule depends only on
$N$, each bid's relation to $m$, independent identically distributed coins,
and the symmetric statistic $K$, proving weak anonymity.

For exact implementation, compute
\[
P_N=2^{-N}\sum_{k=T}^{N-T}\binom Nk
\]
using the recurrence for consecutive binomial coefficients.  This takes
$O(N)$ arithmetic operations on integers with $O(N)$ bits.  The rationals
$r^N$, $\delta_N$, $P_N$, $A_N$, and $w_N$ have bit length polynomial in
$N$, $\log h$, and the public-description length.  Hence their exact
arithmetic and comparisons take polynomial time.  There are $O(N)$ exact
rational-coin calls, each expected-polynomial-time by
\Cref{lem:rational-coin}.  This proves the resource claim and exactness of all
outputs.
\end{proof}

Feasibility uses only the central-window threshold.  The next lemma uses the
padding comparison to identify the normalized expected revenue needed in all
three incentive arguments.

\begin{lemma}[Revenue under fair eligibility and ineligible padding]
\label{lem:large-revenue}
Assume $h\ge64$.
\begin{enumerate}
  \item If a length-$N$ input has $L\ge h$ independent fair eligibility bits
  and every other eligibility bit fixed to zero, where $N\ge L$, then its
  expected miner revenue is at most $cm$.
  \item If all $N\ge h$ eligibility bits are independent and fair, expected
  miner revenue is exactly $cm=hm/128$.
  \item If a length-$N$ input has $N-1\ge h$ independent fair eligibility
  bits and one additional eligibility bit fixed to either zero or one, then,
  conditional on that fixed bit, expected miner revenue is exactly $cm$.
\end{enumerate}
\end{lemma}

\begin{proof}
In the first case, the eligibility count is $B_L$, and
\Cref{lem:binomial-padding} gives
\[
\mathbb{E}[\mu]=w_NA_NF_{L,N}
\le w_NA_NP_N=cm.
\]
In the second case the central event has probability $P_N$, so equality holds.
In the third case, its probabilities for fixed bit zero and fixed bit one are,
respectively,
$F_{N-1,N}$ and $\Pr[T\le B_{N-1}+1\le N-T]$.  Both equal $P_N$ by
\Cref{lem:binomial-padding}, and multiplication by $w_NA_N$ gives $cm$.
\end{proof}

\subsection{Incentive properties for large \texorpdfstring{$h$}{h}}

We first isolate direct user utility.  The decreasing sequence $\delta_N$
makes its truthful optimum nonincreasing in the total number of identities,
which is precisely the monotonicity needed against false names.

\begin{lemma}[False-name truthfulness from decreasing confirmation potential]
\label{lem:large-uic}
For $h\ge64$,
\hyperref[alg:sybil-potential-mechanism]{\textsc{SybilPotentialMechanism}}
is ex post false-name user incentive compatible for every finite vector of
outsider bids.
\end{lemma}

\begin{proof}
At total submitted length $N$, let $a(b)=\Pr[E(b)=1]$.  The expected direct
utility of a value-$v$ designated identity reported as $b$ is
\begin{align}
u_N(v,b)
&=a(b)(v-w_N)+(1-a(b))\delta_Nv\notag\\
&=\delta_Nv+a(b)(1-\delta_N)(v-m).
\label{eq:direct-utility}
\end{align}
Thus truthful reporting maximizes direct utility: it has eligibility
probability zero when $v<m$, probability one when $v>m$, and at $v=m$ the
coefficient of $a(b)$ vanishes.  Its optimum is
\begin{equation}
U_N(v)=
\begin{cases}
\delta_Nv,&0\le v\le m,\\[1ex]
v-m+\delta_Nm,&v\ge m.
\end{cases}
\label{eq:truthful-utility}
\end{equation}
By \Cref{lem:parameter-bounds}, $U_N(v)$ is nonincreasing in $N$.

Fix $s$ outsider bids.  Truthful single-identity participation has length
$n=s+1$ and utility $U_n(v)\ge0$.  A deviation with no designated identity
has no intrinsic value and makes only nonnegative payments, so its utility is
at most zero.  Otherwise its designated identity and finitely many fakes make
the new length $N\ge n$.  The designated identity contributes at most
$U_N(v)\le U_n(v)$, while all fake payments are nonnegative.  This pointwise
comparison for each randomized choice of reports remains valid after
averaging.
\end{proof}

User truthfulness alone does not control the miner's revenue share.  For a
miner-only deviation, the padding lemma handles the case of no eligible fake,
while one eligible fake already pays enough to cover the maximum revenue
excess.

\begin{lemma}[Miner deviations with finite false names]
\label{lem:large-mic}
For $h\ge64$,
\hyperref[alg:sybil-potential-mechanism]{\textsc{SybilPotentialMechanism}}
is exact Bayesian miner incentive compatible for every $\rho\in[0,1]$ and
every randomized finite fake-bid strategy chosen independently of at least
$h$ honest values.
\end{lemma}

\begin{proof}
Let $L\ge h$ honest users bid truthfully.  Their eligibility bits are
independent and fair, so honest miner utility is $\rho cm$ by
\Cref{lem:large-revenue}.  Fix a submitted fake vector and condition on all
of its eligibility outcomes; these are independent of the honest fair bits.
Let $N$ be the total length and let $f$ be the number of eligible fakes.

If $f=0$, all fake payments vanish, and the first part of
\Cref{lem:large-revenue} bounds conditional expected revenue by $cm$.  If
$f\ge1$, each eligible fake pays $w_N$, while
\eqref{eq:maximum-revenue} bounds revenue pointwise by $cm+w_N$.  Conditional
coalition utility is therefore at most
\[
\rho(cm+w_N)-fw_N
=\rho cm-(f-\rho)w_N
\le\rho cm.
\]
The inequality survives averaging over fake eligibility coins and a randomized
choice of any finite fake vector.  Aborting yields zero, also no more than the
honest baseline.
\end{proof}

It remains to combine both effects for a coalition of miners and one user.
The only new configuration has one eligible designated identity padded solely
by ineligible fakes; the potential drop in \eqref{eq:potential-drop} was chosen
to pay for that configuration.

\begin{lemma}[One-user coalition deviations]
\label{lem:large-scp}
For $h\ge64$,
\hyperref[alg:sybil-potential-mechanism]{\textsc{SybilPotentialMechanism}}
is exact Bayesian $1$-side-contract-proof for every $\rho\in[0,1]$ and every
randomized finite joint deviation by the miners and one user, when at least
$h$ honest outsiders have independent values from $D$.
\end{lemma}

\begin{proof}
Let $L\ge h$ be the number of honest outsiders, let the colluding user's value
be $v\ge0$, and set $n=L+1$.  Under truthful single-identity participation,
the user's expected utility is $U_n(v)$.  Conditional on either eligibility
status of that identity, expected miner revenue is $cm$ by the third part of
\Cref{lem:large-revenue}.  Thus truthful coalition utility is
\begin{equation}
U_n(v)+\rho cm.
\label{eq:coalition-baseline}
\end{equation}

Fix a finite deviating bid vector and condition on every eligibility outcome
of its identities.  Let $N$ be the total length, and let $f$ count eligible
fake identities, excluding the designated identity when one exists.

If $f\ge1$, then $N\ge n$.  Conditional on its eligibility, a designated
identity's direct utility is at most $U_N(v)$: an ineligible one gives
$\delta_Nv$, while an eligible one gives $v-w_N$; comparison with
\eqref{eq:truthful-utility} verifies the claim on both sides of $m$.  With no
designated identity, direct utility is zero.  Thus direct utility is at most
$U_N(v)\le U_n(v)$.  Using \eqref{eq:maximum-revenue}, conditional coalition
utility is at most
\[
U_n(v)-fw_N+\rho(cm+w_N)
\le U_n(v)+\rho cm.
\]

Suppose $f=0$ and no identity is designated.  The coalition has no intrinsic
value, all of its identities are ineligible, and its payments vanish.  The
first part of \Cref{lem:large-revenue}, applied to the $L$ fair honest bits,
bounds expected revenue by $cm$.  This includes complete dropout, for which
$N=L$.

Suppose $f=0$ and a designated identity is ineligible.  Then $N\ge n$, its
direct utility is $\delta_Nv\le U_N(v)\le U_n(v)$, all coalition payments
vanish, and the same padding bound limits expected revenue to $cm$.

Finally suppose $f=0$ and the designated identity is eligible.  Here $N\ge n$.
If $N=n$, there are no fakes; its direct utility $v-w_n$ is at most $U_n(v)$,
and conditional expected revenue is exactly $cm$ by
\Cref{lem:large-revenue}.  If $N>n$, all additional identities are ineligible
and cost zero.  Equation \eqref{eq:truthful-utility} gives
\begin{equation}
U_n(v)-(v-w_N)\ge m(\delta_n-\delta_N).
\label{eq:user-potential-loss}
\end{equation}
For $v\ge m$ this is equality; for $v\le m$, the left side exceeds the right
side by $(1-\delta_n)(m-v)$.  Because $n\ge h$ and $N>n$,
\eqref{eq:potential-drop} yields
\[
U_n(v)-(v-w_N)\ge md_N.
\]
At the same time, \eqref{eq:maximum-revenue} bounds expected revenue by
$cm+md_N$.  Hence deviation utility is at most
\[
U_n(v)-md_N+\rho(cm+md_N)
=U_n(v)+\rho cm-(1-\rho)md_N
\le U_n(v)+\rho cm.
\]

These cases exhaust every conditioned finite deviation.  Averaging over
eligibility coins and randomized strategic choices preserves comparison with
\eqref{eq:coalition-baseline}; aborting gives zero and cannot improve that
nonnegative baseline.
\end{proof}

\begin{proof}[Proof of \Cref{thm:main}]
Define $\Pi_{h,D}$ by
\hyperref[alg:parity-mechanism]{\textsc{ParityMechanism}} for $1\le h<64$
and by
\hyperref[alg:sybil-potential-mechanism]{\textsc{SybilPotentialMechanism}}
for $h\ge64$.  Exact balanced eligibility follows from
\Cref{lem:balanced-median}.  In the first regime, all incentive, pointwise
feasibility, anonymity, complexity, and revenue claims follow from
\Cref{lem:parity-guarantees}.  In the second regime, feasibility, anonymity,
and exact expected-polynomial implementation follow from
\Cref{lem:large-implementation}; the three incentive properties follow from
\Cref{lem:large-uic,lem:large-mic,lem:large-scp}.  Each incentive proof allows
an arbitrary finite strategic bid vector.

For every truthful input length $\ell\ge h$,
\Cref{lem:parity-guarantees,lem:large-revenue} gives
\[
\mathbb{E}_{V\sim D^\ell,\,\Pi_{h,D}}[\mu]
=
\begin{cases}
m_D/2,&1\le h<64,\\[1ex]
h m_D/128,&h\ge64.
\end{cases}
\]
When $h<64$, $m_D/2\ge h m_D/128$.  Taking the infimum over $\ell\ge h$
proves the asserted lower bound.
\end{proof}

\bibliographystyle{alpha}
\bibliography{references}

\end{document}
